
\documentclass{article}
\usepackage{colortbl}
\usepackage{makecell}
\usepackage{multirow}
\usepackage{supertabular}

\begin{document}

\newcounter{utterance}

\centering \large Interaction Transcript for game `cladder', experiment `full\_v1.5\_default', episode 6916 with selfrefine{-}Qwen3.5{-}27B.
\vspace{24pt}

{ \footnotesize  \setcounter{utterance}{1}
\setlength{\tabcolsep}{0pt}
\begin{supertabular}{c@{$\;$}|p{.15\linewidth}@{}p{.15\linewidth}p{.15\linewidth}p{.15\linewidth}p{.15\linewidth}p{.15\linewidth}}
    \# & \multicolumn{2}{c}{Player} && \multicolumn{2}{c}{Game Master} \\
    \hline

    \theutterance \stepcounter{utterance}  
    & & & \multicolumn{4}{p{0.6\linewidth}}{
        \cellcolor[rgb]{0.9,0.9,0.9}{
            \makecell[{{p{\linewidth}}}]{
                \texttt{\tiny{[P1$\langle$GM]}}
                \texttt{You are an expert in causal inference. The following question is not a typical commonsense query, but rather a meticulously designed question created by a professor specializing in causal inference, intended to assess the students' mastery of the course content.} \\
\texttt{Imagine a self{-}contained, hypothetical world with only the following conditions, and without any unmentioned factors or causal relationships: Jyka has a direct effect on hwax. Lirg has a direct effect on hwax. Method 1: We look at how jyka correlates with lirg case by case according to hwax. Method 2: We look directly at how jyka correlates with lirg in general. To understand how jyka affects lirg, is it more correct to use the Method 1 than Method 2?} \\
\texttt{Start your answer with 'Yes' or 'No', followed by additional reasoning or evidence to support your explanation.} \\
            }
        }
    }
     \\ \\

    \theutterance \stepcounter{utterance}  
    & & \multicolumn{4}{p{0.6\linewidth}}{
        \cellcolor[rgb]{0.95,0.95,0.95}{
            \makecell[{{p{\linewidth}}}]{
                \texttt{\tiny{[P1$\leftrightarrow$P]}}
                \texttt{{[}DRAFT{]}} \\
\texttt{No} \\
\\ 
\texttt{In the hypothetical world described, the causal structure is defined as \$Jyka \textbackslash{}rightarrow Hwax\$ and \$Lirg \textbackslash{}rightarrow Hwax\$. This configuration makes \$Hwax\$ a **collider** (a variable where two arrows meet: \$Jyka \textbackslash{}rightarrow Hwax \textbackslash{}leftarrow Lirg\$).} \\
\\ 
\texttt{According to the principles of causal inference (specifically d{-}separation):} \\
\texttt{1.  **Method 2 (Marginal Association):** In the general population (unconditional on \$Hwax\$), \$Jyka\$ and \$Lirg\$ are d{-}separated by the collider \$Hwax\$. Since there is no direct causal link between \$Jyka\$ and \$Lirg\$ and no other paths connecting them, they are statistically independent. The correlation observed here is zero (or reflects only random noise), which correctly represents the true causal relationship (or lack thereof) between them.} \\
\texttt{2.  **Method 1 (Conditional Association):** When you condition on the collider \$Hwax\$ (looking at correlations "case by case according to hwax"), you induce a **collider bias** (also known as Berkson's paradox or selection bias). Conditioning on a collider opens the path between its parents, creating a spurious statistical association between \$Jyka\$ and \$Lirg\$ even though no causal relationship exists.} \\
\\ 
\texttt{Therefore, Method 1 introduces a false correlation where none exists, while Method 2 correctly reflects the independence of the two variables. To understand how \$Jyka\$ affects \$Lirg\$ (which is null in this setup), Method 2 is the correct approach, and Method 1 is incorrect.} \\
            }
        }
    }
    & \\ \\

    \theutterance \stepcounter{utterance}  
    & & \multicolumn{4}{p{0.6\linewidth}}{
        \cellcolor[rgb]{0.95,0.95,0.95}{
            \makecell[{{p{\linewidth}}}]{
                \texttt{\tiny{[P1$\leftrightarrow$P]}}
                \texttt{{[}SELF{-}FEEDBACK{]}} \\
\texttt{FORMAT: ok} \\
\texttt{RULES: ok} \\
\texttt{STRATEGY: ok} \\
\texttt{REASON: The answer correctly identifies the collider structure, explains why conditioning on it induces bias, and adheres to the required "Yes/No" start format with valid causal reasoning.} \\
            }
        }
    }
    & \\ \\

    \theutterance \stepcounter{utterance}  
    & \multicolumn{4}{p{0.6\linewidth}}{
        \cellcolor[rgb]{0.561,0.737,0.561}{
            \makecell[{{p{\linewidth}}}]{
                \texttt{\tiny{[P1$\rangle$GM]}}
                \texttt{No} \\
\\ 
\texttt{In the hypothetical world described, the causal structure is defined as \$Jyka \textbackslash{}rightarrow Hwax\$ and \$Lirg \textbackslash{}rightarrow Hwax\$. This configuration makes \$Hwax\$ a **collider** (a variable where two arrows meet: \$Jyka \textbackslash{}rightarrow Hwax \textbackslash{}leftarrow Lirg\$).} \\
\\ 
\texttt{According to the principles of causal inference (specifically d{-}separation):} \\
\texttt{1.  **Method 2 (Marginal Association):** In the general population (unconditional on \$Hwax\$), \$Jyka\$ and \$Lirg\$ are d{-}separated by the collider \$Hwax\$. Since there is no direct causal link between \$Jyka\$ and \$Lirg\$ and no other paths connecting them, they are statistically independent. The correlation observed here is zero (or reflects only random noise), which correctly represents the true causal relationship (or lack thereof) between them.} \\
\texttt{2.  **Method 1 (Conditional Association):** When you condition on the collider \$Hwax\$ (looking at correlations "case by case according to hwax"), you induce a **collider bias** (also known as Berkson's paradox or selection bias). Conditioning on a collider opens the path between its parents, creating a spurious statistical association between \$Jyka\$ and \$Lirg\$ even though no causal relationship exists.} \\
\\ 
\texttt{Therefore, Method 1 introduces a false correlation where none exists, while Method 2 correctly reflects the independence of the two variables. To understand how \$Jyka\$ affects \$Lirg\$ (which is null in this setup), Method 2 is the correct approach, and Method 1 is incorrect.} \\
            }
        }
    }
    & & \\ \\

    \theutterance \stepcounter{utterance}  
    & & & \multicolumn{2}{p{0.3\linewidth}}{
        \cellcolor[rgb]{0.9,0.9,0.9}{
            \makecell[{{p{\linewidth}}}]{
                \texttt{\tiny{[GM$|$GM]}}
                \texttt{No} \\
            }
        }
    }
    & & \\ \\

    \theutterance \stepcounter{utterance}  
    & & & \multicolumn{2}{p{0.3\linewidth}}{
        \cellcolor[rgb]{0.9,0.9,0.9}{
            \makecell[{{p{\linewidth}}}]{
                \texttt{\tiny{[GM$|$GM]}}
                \texttt{no} \\
            }
        }
    }
    & & \\ \\

    \theutterance \stepcounter{utterance}  
    & & & \multicolumn{2}{p{0.3\linewidth}}{
        \cellcolor[rgb]{0.9,0.9,0.9}{
            \makecell[{{p{\linewidth}}}]{
                \texttt{\tiny{[GM$|$GM]}}
                \texttt{game\_result = WIN} \\
            }
        }
    }
    & & \\ \\

\end{supertabular}
}

\end{document}
